% !TeX TXS-program:compile = txs:///pdflatex/[-shell-escape]
% !TeX TXS-program:pdflatex = pdflatex -synctex=1 -interaction=nonstopmode --shell-escape %.tex
\section{Numbers as Strings}\label{sec:11numbers_as_strings}
\subsection{Video - Assembly}
Este es la referencia a un video pasado.
\subsection{A Single Digit}\label{subsec:11-01a_single_digit}

Las computadoras, reciben un montón de input en formato de texto, por ejemplo
cuando pasamos un valor por parámetro como \textbf{\textit{12345}}, este no se recibe como número
sino como 5 bits en formato \textit{ASCII}, y para poder convertir un texto a número
se suele usar la función \mintinline{C}{atoi}.

En el código \textit{ASCII}, el valor de \( 0 \) es \( 0x30 \), \( 1 \) es \( 0x31 \) \dots \( 9 \) es \( 0x39 \)
\begin{codebox}[ASM]{SOlution}
	.intel_syntax noprefix
	.global atoi_digit
	atoi_digit:
	movzx rax, BYTE PTR [rdi]
	sub rax, 0x30
	ret
\end{codebox}
https://youtube.com/shorts/ADbXNFYlePI?si=Lng4NcVHgYMD_Jge
\begin{flagbox}{Flag Capturada}
	pwn.college{sK0lkWe8HVPu1pm3kLYL05j9bHo.0lM4kDNywSMycjN4EzW}
\end{flagbox}
\subsection{Two Digits}\label{subsec:11-02two_digits}

Bueno ahora con 2 dígitos.

\begin{codebox}[ASM]{Solucion}
	.intel_syntax noprefix
	.global atoi_digit
	.global atoi

	atoi_digit:
	movzx rax, BYTE PTR [rdi]
	sub rax, 0x30
	ret

	atoi:
	push rdi

	call atoi_digit
	imul rax, 10

	pop rdi
	push rax

	lea rdi, [rdi + 1]
	call atoi_digit

	pop rcx
	add rax, rcx

	ret
\end{codebox}

\begin{flagbox}{Flag Capturada}
	pwn.college{wo_I1U54BVUXOJFPIR19DZ84Dsl.01M4kDNywSMycjN4EzW}
\end{flagbox}
\subsection{String to Integer}\label{subsec:11-03string_to_integer}
\begin{codebox}[ASM]{Solucion}
	.intel_syntax noprefix
	.global atoi

	atoi:
	xor rax, rax

	loop:
	movzx rcx, byte ptr [rdi]
	cmp rcx, 0
	je done

	sub rcx, 0x30

	imul rax, 10
	add rax, rcx
	inc rdi
	jmp loop

	done:
	ret
\end{codebox}
\begin{flagbox}{Flag Capturada}
	pwn.college{MPBk374KBZ9ZuG8Lf_z-Jifc1Tw.0FN4kDNywSMycjN4EzW}
\end{flagbox}
\subsection{Negative Numbers}\label{subsec:11-04negative_numbers}

Ahora con números negativos.

\begin{codebox}[ASM]{Solucion}
	.intel_syntax noprefix
	.global atoi

	atoi:
	xor rax, rax
	xor r8, r8

	loop:
	movzx rcx, byte ptr [rdi]
	cmp rcx, 0x2d
	je negative
	cmp rcx, 0
	je done

	sub rcx, 0x30

	imul rax, 10
	add rax, rcx
	inc rdi
	jmp loop

	negative:
	mov r8, 1
	inc rdi
	jmp loop
	done:
	cmp r8, 1
	je multi_negative
	ret

	multi_negative:
	imul rax, -1
	ret
\end{codebox}

\begin{flagbox}{Flag Capturada}
	pwn.college{QMhD2L9sERZdiJ8Tm39hK4_FSYA.0VN4kDNywSMycjN4EzW}
\end{flagbox}

\subsection{Where the Numbers Ends}\label{subsec:11-05where_the_numbers_ends}
\subsection{A Whole Program}\label{subsec:11-06a_whole_program}
\subsection{A Single Digit, Back to Text}\label{subsec:11-07a_single_digit_back_to_text}
\subsection{Divide and Remainder}\label{subsec:11-08divide_and_remainder}
\subsection{Drop the Leading Zero}\label{subsec:11-09drop_the_leading_zero}
\subsection{Integer to String}\label{subsec:11-10integer_to_string}
\subsection{Negative Numbers, Back to text}\label{subsec:11-11negative_numbers_back_to_text}
\subsection{Sum Them All}\label{subsec:11-12sum_them_all}
